\documentclass[10pt,a4paper]{amsart}
\usepackage[margin=19mm]{geometry}
\usepackage{amsmath,amssymb,mathtools}
\usepackage[scr=rsfs,cal=euler]{mathalfa}
\usepackage{xfrac}
\usepackage[unicode,hidelinks,bookmarks=false]{hyperref}
\newcommand{\ie}{i.e.\ }
\newcommand{\cf}{cf.\ }
\newcommand{\resp}{respectively\ }
\newcommand{\Subsection}{\subsection}
\providecommand{\nobreakdash}{\nobreak\hbox{-}}
\setlength{\emergencystretch}{2em}
\multlinegap0pt
\usepackage{amssymb}
\usepackage{tikz}
\newtheorem{lemma}{Lemma}
\newcommand{\inv}{{\operatorname{inv}}}
\newcommand{\todo}[1]{\vspace{5 mm}\par \noindent
\marginpar{\textsc{ToDo}} \framebox{\begin{minipage}[c]{0.95
\textwidth}
 #1 \end{minipage}}\vspace{5 mm}\par}
\title{Conjugating full cycles by adjacent transpositions: diameter and sorting time}
\author{Ron M. Adin, Eli Bagno, Yuval Roichman}
\date{Source: arXiv:2601.12597v2 (CC BY 4.0)}
\begin{document}
\maketitle

\noindent\emph{Source and licence.} Conjugating full cycles by adjacent transpositions: diameter and sorting time. Version arXiv:2601.12597v2 (CC BY 4.0), distributed by arXiv under the Creative Commons Attribution 4.0 International licence, http://creativecommons.org/licenses/by/4.0/, as recorded in the arXiv licence metadata for this exact version and shown on the abstract page https://arxiv.org/abs/2601.12597v2. Full licence notice: \texttt{licenses/arxiv-2601.12597-CC-BY-4.0.txt}.\par\smallskip

\noindent\textit{Editorial scope.} Counted unit: the article's lemma obs:4 (source0.tex lines 1343--1354). The article's complete original proof of it is delivered with it (source0.tex lines 1356--1402, one \texttt{proof} environment of 47 lines). No mathematics is written, repaired or re-derived: the statement and the proof are the article's own lines, in the article's own notation, and no retained line is substituted or re-flowed.  No further source-local context is needed: every cross-reference of every retained line resolves inside the two delivered spans.  Nothing was omitted from any retained span, so the package carries no omission record.  No retained line contains a citation, so the wrapper carries no bibliography and no bibliography entry.  No retained line contains a figure environment, an image-inclusion command or drawing code, so no image is delivered and the package declares no asset dependency.  Editorial formatting only, in addition: an AMS \texttt{amsart} preamble that restates the article's own declarations and macros used by the retained lines, namely the environments \texttt{lemma} on separate counters, together with the standard packages those lines need.  The retained environments print in this document as Lemma 1 in order of appearance, whereas the article prints the same items with the numbering of its own full document.  The complete native source of the article as distributed in the arXiv e-print for this exact version is bundled unchanged as \texttt{sources/192962/source0.tex}.
\begin{lemma}\label{obs:4}
For every permutation $\tau\in S_n$ and $0\le k\le n$, 
%satisfies
%\[
%\sum\limits_{j=1}^k \tau(k)\le \binom{n+1}{2}-\binom{n-k+1}{2} 
%\]
%then
\[
    \inv(\tau)
    \le \binom{n-k}{2}-k + \sum_{i=1}^k \tau(i).  
\]
\end{lemma}

\begin{proof} 
Fix $0\le k\le n$, and evaluate 
\[
   %\sum\limits_{i\le k<j} 
   |\{(i,j) :\, i\le k<j,\, \tau(i)>\tau(j)\}|. 
\]
We can assume, without loss of generality, that $\tau(1)<\tau(2)<\cdots<\tau(k)$. 
Then, for every $1\le i\le k$,
\[
    |\{j :\, k<j,\, \tau(i)>\tau(j)\}|=\tau(i)-i.
\]

%\todo{Eli: I suggest changing the last line to : $\tau(i)-i-(i-1)=\tau(i)-i$ as a hint to the calculation and/or adding an example}

Thus  
\[
   %\sum\limits_{i\le k<j} 
   |\{(i,j) :\, i\le k<j,\, \tau(i)>\tau(j)\}| 
   = \sum_{i=1}^k (\tau(i)-i)
   = \sum_{i=1}^k \tau(i) - \binom{k+1}{2}.
\]
Note that, obviously, 
\[
    %\sum\limits_{1\le i<j\le k} 
    |\{(i,j) :\, i< j \le k,\, \tau(i)>\tau(j)\}| +
    %\sum\limits_{k< i<j\le n} 
    |\{(i,j) :\, k < i < j,\, \tau(i)>\tau(j)\}|
    \le \binom{k}{2} + \binom{n-k}{2}. 
\]
Altogether, 
%one concludes that, 
for every fixed $0\le k\le n$,
\begin{align*}
    \inv(\tau)
    = |\{(i,j) :\, i<j,\, \tau(i)>\tau(j)\}| 
    &\le \sum_{i=1}^k \tau(i) - \binom{k+1}{2} + \binom{k}{2} + \binom{n-k}{2} \\
    &= \sum_{i=1}^k \tau(i) - k + \binom{n-k}{2}.
    \qedhere
\end{align*}
%Proof is completed. 
%    \[
%   \inv(\tau)=
%        \sum\limits_{1\le i<j\le k} |\{i<j:\ \tau(i)>\tau(j)\}|+
%    \sum\limits_{i\le k<j} |\{i<j:\ \tau(i)>\tau(j)\}|+
%    \sum\limits_{k< i<j\le n} |\{i<j:\ \tau(i)>\tau(j)\}|.
%    \]
\end{proof}

\end{document}
